\subsection{Forced response in a small-vessel tank test}

We analyzed three regular head-wave runs of the 1.40 m Codevintec CK-14e at 0.611--0.860 Hz \cite{Palma2025MooredASV,Palma2025CK14Data}. We fitted heave and pitch relative to the surge-corrected center-of-gravity wave elevation on the first half of each run. On the held-out half, prediction errors were 15--31\% of the detrended pitch root-mean-square amplitude and 21--27\% for heave. The measured pitch lag ranged from $32^\circ$ to $80^\circ$, and the fundamental pitch-rate amplitude ranged from $17$ to $42^\circ$/s.

We replayed these heave and pitch responses in the 100 m radar geometry with zero surge, roll, and yaw. At 0.860 Hz and $45^\circ$ bearing, the 95th-percentile midpoint error rose from 0.015 m with the prescribed response to 0.296 m with the measured response. All echoes cleared the power gate and all fixes passed at 0.50 m. At 0.25 m, either leveling or spoke timing recovered the same 28.13 percentage points of visits, and leveling after spoke timing recovered zero. Fitting the held-out half instead yielded 36.69 points of overlap. The tank measurements constrain motion. The radar outcomes come from the simulation.
